\chapter{Channel (Platform Adapter)}

\section{Overview}

A Channel is the adapter between OpenClaw and a chat platform. Each channel is responsible for:

\begin{itemize}
  \item Receiving inbound platform messages and converting them to a unified format
  \item Sending agent replies back to the corresponding platform
  \item Handling platform-specific features (attachments, reactions, threads, etc.)
\end{itemize}

Config path: \texttt{channels.<name>.*}

\section{Built-in Channels}

\begin{center}
\begin{tabular}{lll}
\toprule
\textbf{Channel ID} & \textbf{Platform} & \textbf{Login Method} \\
\midrule
\texttt{whatsapp} & WhatsApp & QR code pairing \\
\texttt{telegram} & Telegram & Bot Token \\
\texttt{discord} & Discord & Bot Token \\
\texttt{slack} & Slack & OAuth/Token \\
\texttt{signal} & Signal & Pairing \\
\texttt{imessage} & iMessage (legacy) & macOS \\
\texttt{bluebubbles} & iMessage (recommended) & BlueBubbles server \\
\texttt{webchat} & Web Chat & Built-in \\
\texttt{googlechat} & Google Chat & Service Account \\
\bottomrule
\end{tabular}
\end{center}

Plugins can add more channels (e.g., Mattermost, Zalo, WeChat, etc.).

\section{Channel Login}

\begin{lstlisting}[language=bash]
# Interactive login (all channels)
openclaw channels login

# Channel-specific
openclaw channels login --channel telegram
openclaw channels login --channel discord

# Check status
openclaw channels status --probe
\end{lstlisting}

\section{Access Control}

Each channel supports \textbf{DM policy} and \textbf{group policy}:

\subsection{DM Policy (dmPolicy)}

\begin{itemize}
  \item \texttt{pairing} (default): Unknown senders receive a pairing code; communication requires approval
  \item \texttt{allowlist}: Only users on the allowlist
  \item \texttt{open}: Anyone can DM
  \item \texttt{disabled}: DMs disabled
\end{itemize}

\subsection{Group Policy (groupPolicy)}

\begin{itemize}
  \item \texttt{allowlist} (default): Only groups on the allowlist
  \item \texttt{open}: All groups
  \item \texttt{disabled}: Groups disabled
\end{itemize}

\subsection{Configuration Example}

\begin{lstlisting}[language=json]
{
  channels: {
    whatsapp: {
      dmPolicy: "pairing",
      allowFrom: ["+15555550123"],
      groups: {
        "*": { requireMention: true }
      }
    },
    telegram: {
      enabled: true,
      dmPolicy: "allowlist",
      allowFrom: ["123456789"]
    }
  }
}
\end{lstlisting}

\section{Routing Mechanism}

OpenClaw routes replies \textbf{back to the originating channel}. Routing is deterministic:

\begin{itemize}
  \item The model does \textbf{not} choose the channel
  \item Replies always return to the source of the inbound message
  \item In multi-channel scenarios, Binding rules route to different Agents
\end{itemize}

\subsection{Session Key Format}

\begin{itemize}
  \item DM: \texttt{agent:<agentId>:<mainKey>}
  \item Group: \texttt{agent:<agentId>:<channel>:group:<groupId>}
  \item Thread: Append \texttt{:thread:<threadId>} to the base key
\end{itemize}

\section{Multiple Accounts}

Channels that support multiple accounts (e.g., WhatsApp) use \texttt{accountId} to identify each login instance.
Different \texttt{accountId} values can route to different Agents, enabling a single server to host multiple numbers.
